\scanpage{9}
\mspage{(7)}
\uline{Ces structures sont héritées par les}
\[
\boxed{\begin{aligned}
C^{\bullet\bullet}_{\mathrm{DRpd}}(\mathcal X_*,k)&\overset{\mathrm{dfn}}=C^\bullet_{\mathrm{DRpd}}(\mathcal X_*;k\{T\},k\{T\}^+,T)\\
&=\operatorname{Hom}(\mathcal X_*,C^{\bullet\bullet}_{\mathrm{DRpd}\,*}(k))
\end{aligned}}
\]
et dépendent de façon contravariante de $\mathcal X_*$, (covariante de $k$).

$\Phi_{k*}$ est \deleted{\unsure{?}} \inserted{un $k$-Module $\frac12$ simplicial} homotope à 0 donc $\Gamma^p_k(\Phi_{k*})\otimes\bigwedge^q\Psi_{k*}$ \deleted{\unsure{est}} est homotope à $\Gamma^p_k(0)\otimes\bigwedge^q\Psi_{k*}$, donc hom. à 0 si $p\ne0$.

Donc $C^{p,q}_{\mathrm{DRpd}\,*}(k)$ est hom. à 0 (donc \deleted{flasque} \inserted{acyclique}) si $p=0$. En degré total donné \inserted{$n$}, on a
\diagram{\begin{tikzcd}[column sep=small]
0\arrow[r] & C_*^{n,0}\arrow[r] & C_*^{n-1,1}\arrow[r] & C_*^{n-2,2}\arrow[r,dash] & \cdots & C_*^{\unsure{1},n-1} & C_*^{0,n}\arrow[l]\arrow[r] & 0\arrow[d]\\
& k_*\arrow[u] & & & & & &0\arrow[d]\\
& & & & & & &0
\end{tikzcd}}
i.e. on trouve une résolution \deleted{flasque} \inserted{de long. $n$} de $k_*$ par des \deleted{groupes} \inserted{$k$-Modules} $\frac12$ simpliciaux qui sont acycliques sauf le dernier -- donc on peut la considérer comme une tronquée en l'\deleted{indice} \inserted{degré} $n$ d'une résolution flasque de $k_*$ -- la cohomologie de ses sections sur un $\mathcal X_*$ est donc la coh. de $\mathcal X_*$ tronquée en degré $n$ :
\[
H^{p,q}_{\mathrm{DRpd}}(\mathcal X_*,k)\simeq
\begin{cases}
H^q(\mathcal X_*,k)&\text{si }q\leq p+q\ \text{i.e. }p\geq0\\
0&\text{si }p<0
\end{cases}
\]
\mspage{8}
